% id: 2019-timbo-1-7 | título: Pressão de saturação | fonte: listas/2019/Timbo/lista2019timbo1.pdf, p. 3
Deduza a equação de Clausius-Clapeyron
\[ \frac{dP_s}{dT} = \frac{\mu\lambda}{RT^2}P_s \]
Sendo o calor latente de vaporização $\lambda$, temperatura $T$, pressão de saturação $P_s$ e massa molar $\mu$.

\textit{Dica}:Trabalhe com um ciclo de Carnot infinitesimal em que o trabalho é realizado pelo vapor de água e que ambas as fontes quente e fria são feitas de água, com temperaturas $T_0$ e $T_1$, respectivamente.
